\documentclass{beamer}
% COMSC-3013 Computer Architecture shared Beamer preamble.
% Week decks live one directory below weeks/ and input this file as:
%   % COMSC-3013 Computer Architecture shared Beamer preamble.
% Week decks live one directory below weeks/ and input this file as:
%   % COMSC-3013 Computer Architecture shared Beamer preamble.
% Week decks live one directory below weeks/ and input this file as:
%   \input{../_shared/beamer-preamble.tex}
% Keep the course self-contained. Change visual/build conventions here,
% not by forking one-off week themes.

\usetheme{Boadilla}
\usecolortheme{seahorse}
\setbeamertemplate{navigation symbols}{}
\setbeamertemplate{footline}[frame number]

\usepackage{amsmath}
\usepackage{booktabs}
\usepackage{graphicx}
\usepackage{listings}
\usepackage{pgfpages}

\lstset{
  basicstyle=\ttfamily\small,
  columns=fullflexible,
  keepspaces=true,
  showstringspaces=false,
  breaklines=true
}

% Speaker notes remain hidden in the ordinary student build. Define \shownotes
% before inputting the deck to create an instructor copy; see weeks/README.md.
\ifdefined\shownotes
  \setbeameroption{show notes on second screen=right}
\fi

\newcommand{\ArchTitleSlide}[4]{%
  % #1 week label, #2 title, #3 central question, #4 Wednesday prediction
  \begin{frame}
    \titlepage
  \end{frame}
  \begin{frame}{The machine question}
    \Large #3

    \vspace{1.2em}
    \normalsize\textbf{Prediction to carry into Wednesday:} #4
  \end{frame}
}

% Keep the course self-contained. Change visual/build conventions here,
% not by forking one-off week themes.

\usetheme{Boadilla}
\usecolortheme{seahorse}
\setbeamertemplate{navigation symbols}{}
\setbeamertemplate{footline}[frame number]

\usepackage{amsmath}
\usepackage{booktabs}
\usepackage{graphicx}
\usepackage{listings}
\usepackage{pgfpages}

\lstset{
  basicstyle=\ttfamily\small,
  columns=fullflexible,
  keepspaces=true,
  showstringspaces=false,
  breaklines=true
}

% Speaker notes remain hidden in the ordinary student build. Define \shownotes
% before inputting the deck to create an instructor copy; see weeks/README.md.
\ifdefined\shownotes
  \setbeameroption{show notes on second screen=right}
\fi

\newcommand{\ArchTitleSlide}[4]{%
  % #1 week label, #2 title, #3 central question, #4 Wednesday prediction
  \begin{frame}
    \titlepage
  \end{frame}
  \begin{frame}{The machine question}
    \Large #3

    \vspace{1.2em}
    \normalsize\textbf{Prediction to carry into Wednesday:} #4
  \end{frame}
}

% Keep the course self-contained. Change visual/build conventions here,
% not by forking one-off week themes.

\usetheme{Boadilla}
\usecolortheme{seahorse}
\setbeamertemplate{navigation symbols}{}
\setbeamertemplate{footline}[frame number]

\usepackage{amsmath}
\usepackage{booktabs}
\usepackage{graphicx}
\usepackage{listings}
\usepackage{pgfpages}

\lstset{
  basicstyle=\ttfamily\small,
  columns=fullflexible,
  keepspaces=true,
  showstringspaces=false,
  breaklines=true
}

% Speaker notes remain hidden in the ordinary student build. Define \shownotes
% before inputting the deck to create an instructor copy; see weeks/README.md.
\ifdefined\shownotes
  \setbeameroption{show notes on second screen=right}
\fi

\newcommand{\ArchTitleSlide}[4]{%
  % #1 week label, #2 title, #3 central question, #4 Wednesday prediction
  \begin{frame}
    \titlepage
  \end{frame}
  \begin{frame}{The machine question}
    \Large #3

    \vspace{1.2em}
    \normalsize\textbf{Prediction to carry into Wednesday:} #4
  \end{frame}
}


\title{Week NN Monday}
\subtitle{WEEK TITLE}
\author{COMSC-3013 Computer Architecture}
\date{}

\begin{document}

\begin{frame}
  \titlepage
\end{frame}

\ArchQuestionSlide{REPLACE CENTRAL QUESTION}{REPLACE BOUNDED PREDICTION}
\note{Open by naming the prior belief this week will test. Do not begin with a vocabulary dump.}

\begin{frame}{The useful model}
\begin{itemize}
  \item REPLACE: smallest mechanism needed for Wednesday.
  \item REPLACE: relationship the student should be able to predict from.
  \item REPLACE: scope or assumption that matters.
\end{itemize}
\note{Keep the model small. If a concept does not help interpret Wednesday's evidence or Week 14 judgment, consider omitting it.}
\end{frame}

\begin{frame}[fragile]{Persistent worked example}
\begin{lstlisting}
REPLACE WITH TRACE, CODE, DISASSEMBLY, EQUATION, OR PSEUDOCODE
\end{lstlisting}
\note{Reuse a familiar specimen or workload whenever possible. Work it far enough that students recognize the same object in the lab.}
\end{frame}

\begin{frame}{What this model does not prove}
\begin{itemize}
  \item REPLACE WITH ONE IMPORTANT LIMITATION.
  \item REPLACE WITH ONE LIKELY OVERCLAIM TO AVOID.
\end{itemize}
\note{Scope is part of Architecture reasoning, not a disclaimer pasted on at the end.}
\end{frame}

\begin{frame}{Wednesday: make the machine argue}
\Large
If our model is right, changing \textbf{X} while holding \textbf{Y} reasonably stable should cause \textbf{Z}.

\vspace{1em}
\normalsize
Wednesday we measure whether the machine agrees.
\note{End before resolving the experiment. Students should carry a real prediction into Wednesday.}
\end{frame}

\end{document}
